\chapter{三阶AKNS方程组的逆时空约化}
现在我们将眼光转向 $ 3 \times 3 $ 逆时空非局域 AKNS 系统, 读者在学完本章节后会更加体会到如何利用 RH 方法处理反散射数据的对称性求解逆时空问题.

本章接分为四个部分, 首先介绍三阶 AKNS 方程组的约化, 这一部分基本照抄了君言师姐的论文; 第二部分为利用 RH 方法求解耦合 NLS 方程的 $N$ 孤子解, 这部分参考了王灯山的论文\cite{Wang2010JMP}; 第三部分则证明了特征值问题对称性与 RH 问题对称性的等价性, 这一部分用以回答上一章节问题 \ref{ch3-question-RH-and-eigenvalue-problem}: 为何特征值和特征向量的对称性可替换 $N$ 孤子解相关参数? 这部分为笔者的原创内容; 第四部分则介绍利用 RH 方法求解三阶逆时空 NLS 方程组的 N 孤子解, 构成了笔者的第一篇文章.

\section{AKNS 方程组的约化}
首先考虑 $ 3 \times 3 $ AKNS 系统, 它的 Lax 对为
\begin{align}
    \Phi_{x} & = U \Phi = (-\rmi \zeta \Lambda + Q)\Phi, \label{ch4-eq-AKNS-Lax-Pair-x}       \\
    \Phi_{t} & = V \Phi = (2 \rmi \zeta^{2} \Lambda + R)\Phi,  \label{ch4-eq-AKNS-Lax-Pair-t}
\end{align}
其中 $R=- 2\zeta Q + \rmi(Q^{2} + Q_{x})\Lambda $, 且
\begin{equation}
    Q = \begin{pmatrix}
        0      & 0      & p \\
        0      & 0      & q \\
        -r_{1} & -r_{2} & 0
    \end{pmatrix}, \quad
    \Lambda = \begin{pmatrix}
        -1 & 0  & 0 \\
        0  & -1 & 0 \\
        0  & 0  & 1
    \end{pmatrix}
\end{equation}
其中 $ r_{1} = ap_{1} + bq_{1}, r_{2} = cp_{1} + dq_{1}$, 这里 $ p_{1} $ 可能为 $ p, p(-x), p(-t), p(-x,-t)$ 及其共轭, $ q_{1} $同理, 且 $ a, b, c, d \in \mathbb{C} $ 为任意复数. 此时零曲率方程 $U_{t} - V_{x} + [U,V] = 0$ 等价于 $\rmi Q_{t} \Lambda + Q_{xx} - 2 Q^{3} = 0$, 则得到四分量方程
\begin{align}
    (13): \quad & \rmi p_{t} + p_{xx} + 2p(pr_{1} + qr_{2}) = 0,    \label{ch4-eq-4-component-AKNS-components-13}     \\
    (23): \quad & \rmi q_{t} + q_{xx} + 2q(pr_{1} + qr_{2}) = 0,   \label{ch4-eq-4-component-AKNS-components-23}      \\
    (31): \quad & -\rmi r_{1t} + r_{1xx} + 2r_{1}(pr_{1} + qr_{2}) = 0, \label{ch4-eq-4-component-AKNS-components-31} \\
    (32): \quad & -\rmi r_{2t} + r_{2xx} + 2r_{2}(pr_{1} + qr_{2}) = 0. \label{ch4-eq-4-component-AKNS-components-32}
\end{align}
考虑相容性条件. 我们有 $ p_{1} $ 只可能为 $ p^{*}(x,t), p^{*}(-x,t), p(x,-t), p(-x,-t) $, $ q_1 $ 同理. 则得到如下四种约化
\begin{align}
    r_{1} & = a p^{*}(x,t) + b q^{*}(x,t), \quad r_{2} = b^{*} p^{*}(x,t) + d q^{*}(x,t), \quad a, d \in \mathbb{R} \label{ch4-eq-reduction-local},         \\
    r_{1} & = a p^{*}(-x,t) + b q^{*}(-x,t), \quad r_{2} = b^{*} p^{*}(-x,t) + d q^{*}(-x,t), \quad a, d \in \mathbb{R} \label{ch4-eq-reduction-reverse-x}, \\
    r_{1} & = a p(x,-t) + b q(x,-t), \quad r_{2} = b p(x,-t) + d q(x,-t), \label{ch4-eq-reduction-reverse-t}                                                \\
    r_{1} & = a p(-x,-t) + b q(-x,-t) , \quad r_{2} = c p(-x,-t) + d q(-x,-t), \label{ch4-eq-reduction-reverse-x-t}
\end{align}
\begin{remark}
    这里的相容性条件指通过 \eqref{ch4-eq-4-component-AKNS-components-13}-\eqref{ch4-eq-4-component-AKNS-components-32} 划分为两份量方程, 考虑利用 (13), (23) 式化为 (31) 式, 要将 $\rmi p_{t} $ 变为 $ - \rmi r_{1t} $, 则要么 $ t \to -t $, 要么做共轭. 由于对 $ t $ 求导, 故与 $ x $ 无关, 只能选择 上述四种情形.
\end{remark}
若进一步考虑将其约化为单分量方程, 则有如下三种可能
\begin{align}
    q(x,t) & = p(-x,t), \quad a = d, \label{ch4-eq-reduction-q1}          \\
    q(x,t) & = p^{*}(x,-t), \quad a = d^{*} \label{ch4-eq-reduction-q2},  \\
    q(x,t) & = p^{*}(-x,-t), \quad a = d^{*} \label{ch4-eq-reduction-q3},
\end{align}
接下来我们讨论可能的约化情况, 为表示方便, 记 $ \Delta = pr_{1} + qr_{2} $.
\subsection*{情形 1. $p_{1} = p^{*}, q_{1} = q^{*}$}
此时我们有 $ r_1 = ap^{*} + bq^{*}, r_2 = cp^{*} + dq^{*} $, 带入 \eqref{ch4-eq-4-component-AKNS-components-13}-\eqref{ch4-eq-4-component-AKNS-components-32}, 得到
\begin{align}
    (13): & \quad ip_{t} + p_{xx} + 2p \Delta = 0                                                                           \\
    (23): & \quad iq_{t} + q_{xx} + 2q \Delta = 0                                                                           \\
    (31): & \quad a(- \rmi p^{*}_{t} + p_{xx}^{*} + 2 p^{*} \Delta) + b(- \rmi q^{*}_{t} + q_{xx}^{*} + 2 q^{*} \Delta) = 0 \\
    (32): & \quad c(- \rmi p^{*}_{t} + p_{xx}^{*} + 2 p^{*} \Delta) + d(- \rmi q^{*}_{t} + q_{xx}^{*} + 2 q^{*} \Delta) = 0
\end{align}
要利用 (13), (23) 式消去 (31), (32) 式, 则需要 $\Delta = \Delta^{*} $,
\begin{align}
    \Delta     & = a p p^{*} + b p q^{*} + c p^{*} q + d q q^{*}                 \\
    \Delta^{*} & = a^{*} p^{*} p + b^{*} p^{*} q + c^{*} p q^{*} + d^{*} q^{*} q
\end{align}
可得 $ a, d \in \mathbb{R}, c = b^{*} $. 此时我们有
\begin{equation}
    (31) = a (12)^{*} + b (13)^{*} \quad (32) = b^{*} (12)^{*} + d (13)^{*}
\end{equation}
方程约化为
\begin{align}
    \rmi p_{t} + p_{xx} + 2p \left( a p p^{*} + b p q^{*} + b^{*} p^{*} q + d q q^{*} \right) & = 0 \label{ch4-eq-Coupled-NLS-p} \\
    \rmi q_{t} + q_{xx} + 2q \left( a p p^{*} + b p q^{*} + b^{*} p^{*} q + d q q^{*} \right) & = 0 \label{ch4-eq-Coupled-NLS-q}
\end{align}
这个方程即为文献\cite{Wang2010JMP} 提出的具有四波混频效应的耦合广义NLS方程. 下面我们进一步约化, 尝试把两个方程约化为一个方程. 这里 $ q $ 可取 $ p(-x), p^{*}(-t), p^{*}(-x,-t) $.
\subsubsection*{Case 1.1 $ q = p(-x) $}
带入 \eqref{ch4-eq-Coupled-NLS-p}-\eqref{ch4-eq-Coupled-NLS-q}, 得到 $ a=d, b \in \mathbb{R} $, 此时$ r_{1} = ap^{*} + bp(-x), r_{2} = b^{*} p^{*} + ap(-x) $, 方程约化为
\begin{equation}
    \rmi p_{t} + p_{xx} + 2p \left[ a p p^{*} + b p p^{*}(-x) + b p^{*} p(-x) + a p(-x) p^{*}(-x) \right] = 0
\end{equation}
即单分量逆空间非局域 NLS 方程.
\subsubsection*{Case 1.2 $ q = p^{*}(-t) $}
此时 $ a=d $, 则 $ r_{1} = a p^{*} + b p(-x,-t), r_{2} = b^{*} p^{*} + a p(-x,-t) $ 方程约化为
\begin{equation}
    \rmi p_{t} + p_{xx} + 2p \left[ a p p^{*} + b p p(-t) + b^{*} p^{*} p^{*}(-t) + a p(-t) p^{*}(-t) \right] = 0
\end{equation}
即单分量逆时间非局域 NLS 方程.
\subsubsection*{Case 1.3 $ q = p^{*}(-x,-t) $}
此时 $ a=d $, 则 $ r_{1} = a p^{*} + b p(-x,-t), r_{2} = b^{*} p^{*} + a p(-x,-t) $ 方程约化为
\begin{equation}
    \rmi p_{t} + p_{xx} + 2p \left[ a p p^{*} + b p p(-x,-t) + b^{*} p^{*} p^{*}(-x,-t) + a p(-x,-t) p^{*}(-x,-t) \right] = 0
\end{equation}

\subsection*{情形 2. $ p_1 = p^{*}(-x), q_1 = q^{*}(-x) $}
此时 $ r_1 = ap^{*}(-x) + bq^{*}(-x), r_2 = cp^{*}(-x) + dq^{*}(-x) $, 要消去 \eqref{ch4-eq-4-component-AKNS-components-13}-\eqref{ch4-eq-4-component-AKNS-components-32} 中 $(31), (32)$ 式, 只需 $ \Delta = \Delta^{*}(-x) $, 其中
\begin{align}
    \Delta         & = ap^{*}(-x)p + bq^{*}(-x)p + c^{*}p^{*}(-x)q + dq^{*}(-x)q             \\
    \Delta^{*}(-x) & = a^{*}pp^{*}(-x) + b^{*}qp^{*}(-x) + c^{*}pq^{*}(-x) + d^{*}qq^{*}(-x)
\end{align}
则 $ a,d \in \mathbb{R}, b = c^{*} \in \mathbb{C} $, 此时
\begin{equation}
    (31) = a(13)^{*}(-x) + b(23)^{*}(-x), \quad
    (32) = b^{*}(13)(-x) + d(23)^{*}(-x),
\end{equation}
方程 \eqref{ch4-eq-4-component-AKNS-components-13}-\eqref{ch4-eq-4-component-AKNS-components-32} 约化为
\begin{align}
    \rmi p_{t} + p_{xx} + 2p \left[ a p p^{*}(-x) + b p q^{*}(-x) + b^{*}p^{*}(-x)q + d q^{*}(-x) q \right] & = 0, \label{ch4-eq-reduction-inverse-x-p} \\
    \rmi q_{t} + q_{xx} + 2q \left[ a p p^{*}(-x) + b p q^{*}(-x) + b^{*}p^{*}(-x)q + d q^{*}(-x) q \right] & = 0\label{ch4-eq-reduction-inverse-x-q}
\end{align}
则 q 可取 $ p^{*}(-t), p(-x), p^{*}(-x,-t) $.

\subsubsection*{Case 2.1 $ q = p(-x) $}
此时 $ a=d $, 方程约化为
\begin{equation}
    ip_t + p_{xx} + 2p \left[ app^{*}(-x) + bpp^{*} + bp^{*}(-x)p(-x) + ap^{*}p(-x) \right] = 0
\end{equation}

\subsubsection*{Case 2.2 $q = p^{*}(-t)$}
此时 $ a = d $. 方程约化为
\begin{equation}
    ip_t + p_{xx} + 2p \left[ app^{*}(-x) + bpp(-x,-t) + b^{*}p^{*}(-x)p^{*}(-t) + ap(-x,-t)p^{*}(-t) \right] = 0
\end{equation}

\subsubsection*{Case 2.3 $ q = p^{*}(-x,-t) $}
此时 $r_1 = ap^{*}(-x) + bp(-t), r_2 = b^{*}p^{*}(-t) + ap(-t) $, 方程约化为
\begin{equation}
    ip_t + p_{xx} + 2p \left[ app^{*}(-x) + bpp(-t) + b^{*}p^{*}(-x)p^{*}(-x,-t) + ap(-t)p^{*}(-x,-t) \right] = 0
\end{equation}

\subsection*{情形 3. $ p_1 = p(-t), q_1 = q(-t) $}
此时 $r_1 = ap(-t) + bq(-t), r_2 = cp(-t) + dq(-t)$
要约化 (31) (32), 需要 $ \Delta(-t) = \Delta $, 其中
\begin{equation}
    \begin{aligned}
        \Delta     & = app(-t) + bpq(-t) + cp(-t)q + dqq(-t) \\
        \Delta(-t) & = ap(-t)p + bp(-t)q + cpq(-t) + dq(-t)q
    \end{aligned}
\end{equation}
则 $ b=c $, 此时有
\begin{equation}
    (31) = a(13)(-t) + b(23)(-t), \quad
    (32) = b(13)(-t) + d(23)(-t),
\end{equation}
方程约化为
\begin{align}
    ip_{t} + p_{xx} + 2p \left[ app(-t) + bpq(-t) + bp(-t)q + dqq(-t) \right] = 0, \label{ch4-eq-reduction-inverse-t-p} \\
    iq_{t} + q_{xx} + 2q \left[ app(-t) + bpq(-t) + bp(-t)q + dqq(-t) \right] = 0, \label{ch4-eq-reduction-inverse-t-q}
\end{align}
q 取值同上.

\subsubsection*{Case 3.1 $ q = p(-x) $}
此时 $ a = d $, 方程约化为
\begin{equation}
    ip_t + p_{xx} + 2p \left[ app(-t) + bpp(-x,-t) + bp(-t)p(-x) + ap(-x)p(-x,-t) \right] = 0
\end{equation}

\subsubsection*{Case 3.2 $ q = p^{*}(-t) $}
此时 $ a^{*} = d, b \in \mathbb{R} $, 方程约化为
\begin{equation}
    ip_t + p_{xx} + 2p \left[ app(-t) + bpp^{*} + bp(-t)p^{*}(-t) + a^{*}p^{*}p^{*}(-t) \right] = 0
\end{equation}

\subsubsection*{Case 3.3 $ q = p^{*}(-x,-t) $}
此时 $ a^{*} = d , b \in \mathbb{R}$, 方程约化为
\begin{equation}
    ip_t + p_{xx} + 2p \left[ app(-t) + bpp^{*}(-x) + bp(-t)p^{*}(-x,-t) + a^{*}p^{*}(-x)p^{*}(-x,-t) \right] = 0
\end{equation}

\subsection*{情形 4. $ p_{1} = p(-x,-t), q_{1} = q(-x,-t) $}
此时有 $r_1 = ap(-x,-t) + bq(-x,-t), r_2 = cp(-x,-t) + dq(-x,-t) $, 要约化 (31) (32), 需要 $ \Delta(-x,-t) = \Delta $, 其中
\begin{align}
    \Delta        & = app(-x,-t) + bpq(-x,-t) + cp(-x,-t)q + dq(-x,-t)q \\
    \Delta(-x,-t) & = ap(-x,-t)p + bp(-x,-t)q + cpq(-x,-t) + dqq(-x,-t)
\end{align}
则 $ \Delta(-x,-t) = \Delta $, 只需 $ b=c $, 此时有
\begin{align}
    (31) & = a(13)(-x,-t) + b(23)(-x,-t) \\
    (32) & = b(13)(-x,-t) + d(23)(-x,-t)
\end{align}
方程约化为
\begin{align}
    ip_{t} + p_{xx} + 2p \left[ app(-x,-t) + bpq(-x,-t) + bp(-x,-t)q + dqq(-x,-t) \right] = 0, \label{ch4-eq-reduction-inverse-xt-p} \\
    iq_{t} + q_{xx} + 2q \left[ app(-x,-t) + bpq(-x,-t) + bp(-x,-t)q + dqq(-x,-t) \right] = 0, \label{ch4-eq-reduction-inverse-xt-q}
\end{align}

\subsubsection*{Case 4.1 继续增加约化条件 $ q = p(-x,t) $}
此时 $ a = d $, 则 $ r_1 = ap(-x,-t) + bp(-t), r_2 = bp(-x,-t) + dp(-t) $, 此时有 $ (13) = (12)(-x) $. 方程约化为
\begin{equation}
    ip_t + p_{xx} + 2p \big( app(-x,-t) + bpp(-t) + bp(-x,-t)p(-x) + ap(-t)p(-x) \big) = 0
\end{equation}

\subsubsection*{Case 4.2 继续增加约化条件 $ q = p^{*}(-x) $}
此时 $ a^{*} = d , b \in \mathbb{R}$, 则 $ r_1 = ap(-x,-t) + bp^{*}(-x), r_2 = bp(-x,-t) + dp^{*}(-x)$, 方程约化为
\begin{equation}
    ip_t + p_{xx} + 2p \left[ app(-x,-t) + bpp(-t) + bp(-x,-t)p(-x) + ap(-t)p(-x) \right] = 0
\end{equation}

\subsubsection*{Case 4.3 继续增加约化条件 $ q = p^{*}(-x,-t) $}
此时 $ a = d^{*}, b \in \mathbb{R} $, 则 $ r_1 = ap(-x,-t) + bp^{*}, r_2 = bp(-x,-t) + dp^{*} $, 方程约化为
\begin{equation}
    ip_t + p_{xx} + 2p \left[ app(-x,-t) + bpp^{*} + bp(-x,-t)p^{*}(-x,-t) + a^{*}p^{*}p^{*}(-x,-t) \right] = 0
\end{equation}
至此, 我们推导出所有可能的约化情况.

\section{利用 RH 方法求解广义耦合 NLS 方程}
本部分我们将利用 RH 方法求解零边界的耦合 NLS 方程 \eqref{ch4-eq-Coupled-NLS-p}, \eqref{ch4-eq-Coupled-NLS-q}. 故 $ x \to \pm \infty $ 时, $ p, q \to 0 $, 则
\begin{equation}
    \phi \sim E := e^{\theta \Lambda}, \quad \theta = -\rmi \zeta x + 2 \rmi \zeta^{2} t
\end{equation}
注意到
\begin{equation}
    E^{*}(\zeta^{*}) = E^{\dagger}(\zeta^{*}) = E^{-1}, \quad E^{*}(-x, t-\zeta^{*}) = E^{\dagger}(-x, t-\zeta^{*}) = E^{-1}
\end{equation}
Jost 渐近解为 $ J = \phi E^{-1} $, 则有 $J \to I, x \to \pm \infty$. 将其带入 Lax 对 \eqref{ch4-eq-AKNS-Lax-Pair-x} 有
\begin{equation}
    J_{x}E + J E(- \rmi \zeta \Lambda) = (-\rmi \zeta \Lambda + Q) J E \implies J_{x} = - \rmi \zeta [\Lambda, J] + Q J
    \label{ch4-eq-Lax-mod-x}
\end{equation}
同理有
\begin{equation}
    J_{t} = 2 \rmi \zeta^{2} [\Lambda, J] + R J
    \label{ch4-eq-Lax-mod-t}
\end{equation}
则由上式易见
\begin{equation}
    \tr(J^{-1} J_{x}) = - \rmi \zeta (\tr(J^{-1} \Lambda J) - \tr(\Lambda)) + \tr(J^{-1}Q J) = 0, \quad
\end{equation}
同理有 $\tr(J^{-1} J_{t}) = 0$, 由 Abel 定理有 $\det(J)_{x} = \det(J) \tr(J^{-1} J_{x}) = 0, \det(J)_{t} = \det(J) \tr(J^{-1} J_{t}) = 0$, 则 $\det(J) = 1$. 对上式变换有
\begin{equation}
    \begin{aligned}
        \hat{E} J_{x} + \rmi \zeta \hat{E} [\Lambda, J]   & = \hat{E} Q J, \\
        \hat{E} J_{t} - 2 \rmi \zeta \hat{E} [\Lambda, J] & = \hat{E} R J.
    \end{aligned}
\end{equation}
其中 $\hat{E} X = E^{-1} X E$, 则易得全微分形式
\begin{equation}
    \rmd(\hat{E} J) = \hat{E} [(Q \rmd x + R \rmd t) J]. \label{ch4-eq-reverse-x-volterra}
\end{equation}
\subsection{解析性}
利用全微分形式 \eqref{ch4-eq-reverse-x-volterra}, 可得 Volterra 积分形式为
\begin{equation}
    \begin{aligned}
        J_{-}(x, t, \zeta) & = I + \int_{-\infty}^{x} e^{\rmi \zeta (y-x) \Lambda} Q(y, t) J_{-}(y, t, \zeta) e^{-\rmi \zeta (y-x) \Lambda} \, \rmd y, \\
        J_{+}(x, t, \zeta) & = I - \int_{x}^{+\infty} e^{\rmi \zeta (y-x) \Lambda} Q(y, t) J_{+}(y, t, \zeta) e^{-\rmi \zeta (y-x) \Lambda} \, \rmd y.
    \end{aligned}
\end{equation}
\begin{proof}
    这里仅证明 $J_{-}$, $J_{+}$ 同理对 \eqref{ch4-eq-reverse-x-volterra} 沿 $ -\infty \to x $ 积分, 可得
    \begin{equation}
        \begin{aligned}
            \hat{E} J_{-}(x, t, \zeta)|_{-\infty}^{x} & = \int_{-\infty}^{x} \hat{E}(y,t) Q(y, t) J_{-}(y, t, \zeta) \rmd y                                                   \\
            \implies J_{-}                            & = I + \int_{-\infty}^{x} E(x,t) E^{-1}(y, t) Q(y, t) J_{-}(y, t, \zeta) E(y, t) E^{-1}(x, t) \rmd y                   \\
                                                      & = I + \int_{-\infty}^{x} e^{\rmi \zeta (y-x) \Lambda} Q(y, t) J_{-}(y, t, \zeta) e^{-\rmi \zeta (y-x) \Lambda} \rmd y
        \end{aligned}
        \label{ch4-eq-volterra-J}
    \end{equation}
    若设 $J_{-} = (J_{-,1}, J_{-,2}, J_{-,3})$, 现考虑 $J_{-,1}$, 由上式得
    \begin{equation}
        \begin{aligned}
            J_{-,1}(x, t, \zeta) & = (1,0,0)^{T} + \int_{-\infty}^{x} e^{\rmi \zeta (y-x) \Lambda} Q(y, t) J_{-}(y, t, \zeta) e^{\rmi \zeta (y-x)}(1,0,0)^{T} \rmd y \\
                                 & = (1,0,0)^{T} + \int_{-\infty}^{x} D_{1}(x,y,\zeta) Q(y, t) J_{-,1}(y, t, \zeta)\rmd y
        \end{aligned}
    \end{equation}
    其中 $D_{1}(x,y,\zeta) = \diag( 1,1,e^{2\rmi\zeta(y-x)})$.
    定义 Neumann 级数及递推序列
    \begin{equation}
        \begin{aligned}
            J_{-,1}(x,t,\zeta) = \sum_{n=0}^{\infty}c_{n},      \quad
            c_{0} = (1,0,0)^{T}, \quad c_{n+1}  = \int_{-\infty}^{x} D_{1}(x,y,\zeta) Q(y,t)c_{n}(y,t,\zeta) \rmd y.
        \end{aligned}
        \label{ch4-eq-J-minus-Neumann-series}
    \end{equation}
    对向量取 $\ell^{1}$ 范数 $\|v\|_{1}=\sum_{j=1}^{3}|v_{j}|$, 对矩阵取诱导范数 $\|A\|=\max_{1\leq k\leq3}\sum_{j=1}^{3}|A_{jk}|$.若 $\zeta \in \mathbb{C}_{-}$, 有 $\|D(x,y,\zeta)\|\leq 1$, 则
    \begin{equation}
        \|c_{n+1}\| \leq \int_{-\infty}^{x} \|Q(y,t)\| \|c_{n}(y,t,\zeta)\| \rmd y.
    \end{equation}
    记
    \begin{equation}
        \omega = \int_{-\infty}^{x}\|Q(y,t)\|\rmd y,
    \end{equation}
    由 $Q \in L^{1}$, 则任意固定 $x \in \mathbb{R}$, 上述积分存在, 且 $\omega_{x} = \| Q \|$.

    易见 $\| c_{1} \| \leq \omega$, 由数学归纳法, 不妨设 $\| c_{n} \| \leq \omega^{n}/n!$, 则
    \begin{equation}
        \|c_{n+1}\| \leq \int_{-\infty}^{x} \omega_{x} \frac{\omega^{n}}{n!} \rmd y = \frac{\omega^{n+1}}{(n+1)!}.
    \end{equation}
    故递推序列存在且有界, 又由于
    \begin{equation}
        \sum_{n=0}^{\infty} \| c_{n} \| \leq \sum_{n=0}^{\infty} \frac{\omega^{n}}{n!} = e^{\omega}.
    \end{equation}
    由 Weierstrass 判别法, Neumann 级数 \eqref{ch4-eq-J-minus-Neumann-series}在 $\mathbb{C}_{-}$ 上一致收敛, 故 $J_{-,1}(x,t,\zeta)$ 在 $\mathbb{C}_{-}$中解析且连续到实轴.

    唯一性: 设$J'_{-,1}$ 也为 Volterra 积分的解，令 $h = J_{-,1} - J'_{-,1}$
    则 $h$ 也为 Volterra 积分方程的解，即
    \begin{equation}
        h(x,t,\zeta) = \int_{-\infty}^{x} D_{1}(x,y,\zeta) Q(y,t)h(y,t,\zeta)\,\rmd y,
        \label{ch4-eq-H-minus-first}
    \end{equation}
    由 Grönwall 不等式, $h \equiv 0$, 唯一性得证.

    同理可得$J_{-,1}, J_{-,2}, J_{+,3}$ 在 $\mathbb{C}_{-}$ 中解析且连续到边界., $J_{+,1}, J_{+,2}, J_{-,3}$ 在 $\mathbb{C}_{+}$ 解析且连续到边界.
\end{proof}
\subsection{对称性}
为表示方便, 我们记 $\phi_{\pm} = J_{\pm} E$, 则 $\phi_{\pm}$ 为 Lax 对 \eqref{ch4-eq-AKNS-Lax-Pair-x} 的解, 且满足对称关系
\begin{equation}
    \phi_{-} = \phi_{+} S(\zeta) \implies J_{-} E = J_{+} E S(\zeta), \quad \zeta \in \mathbb{R},
\end{equation}
则由于 $\det(J_{\pm}) = 1$, 则 $\det(S(\zeta)) = 1$. 若记
\begin{equation}
    P^{+} = [\phi_{+,1}, \phi_{+,2}, \phi_{-,3}] E^{-1} = J_{+}H_{2} + J_{-}H_{1}, \quad \zeta \in \mathbb{C}_{+} \label{ch4-eq-RHP-P1}
\end{equation}
其中 $ H_{1} = \diag(0,0,1), H_{2} = \diag(1,1,0) $, 则 $ P^{+} $ 在 $ \mathbb{C}_{+} $ 解析, 联系 \eqref{ch4-eq-RHP-P1} 及 Volterra 积分\eqref{ch4-eq-volterra-J}, 可得
\begin{equation}
    P^{+}(x,\zeta) \to I,  (\zeta \in \mathbb{C}_{+}),
\end{equation}
我们现在考虑 \eqref{ch4-eq-Lax-mod-x} 的伴随问题
\begin{equation}
    K_{x} = - \rmi\zeta [\Lambda, K] - K Q, \label{ch4-eq-Lax-adjoint}
\end{equation}
并设 $ J_{\pm}^{-1} $ 为伴随问题的 Jost 矩阵, 若记
\begin{equation}
    \phi_{\pm}^{-1} = (\hat{\phi}_{\pm,1}, \hat{\phi}_{\pm,2}, \hat{\phi}_{\pm,3})^{T},
\end{equation}
重复上述步骤有
\begin{equation}
    \hat{\phi}_{+,1}, \hat{\phi}_{+,2}, \hat{\phi}_{-,3} \in \mathbb{C}_{-}, \quad \hat{\phi}_{-,1}, \hat{\phi}_{-,2}, \hat{\phi}_{+,3} \in \mathbb{C}_{+}
\end{equation}
可得
\begin{equation}
    P^{-} = E (\hat{\phi}_{-,1}, \hat{\phi}_{-,2}, \hat{\phi}_{+,3})^{T} = H_{2}J_{+}^{-1} + H_{1} J_{-}^{-1}, \quad \zeta \in \mathbb{C}_{-} \label{ch4-eq-RHP-P-minus}
\end{equation}
同理有 $P^{-}(x,\zeta) \to I(\zeta \in \mathbb{C}_{-})$. 接下来我们考虑对称性, 为表示方便记
\begin{equation}
    M_{0} = \diag(A_{1}, 1), \quad M_{1} = \diag(A_{1}, -1), \quad M_{2} = \diag(A_{2}, 1), \quad M_{3} = \diag(A_{2}, -1)
\end{equation}
其中
\begin{equation}A_{1} = \begin{pmatrix}
        a & b^{*} \\
        b & d
    \end{pmatrix}, \quad A_{2} = \begin{pmatrix}
        a & b \\
        b & d
    \end{pmatrix}.
\end{equation}
容易发现 $ Q^{\dagger} = - M_{0} Q M_{0}^{-1}$, 则
\begin{proposition}
    耦合 NLS 满足如下对称关系
    \begin{equation}
        J_{\pm}^{\dagger}(\zeta^{*}) M_{0} = M_{0} J_{\pm}^{-1}(\zeta), \quad S^{\dagger}(\zeta^{*}) M_{0} = M_{0} S^{-1} . \label{ch4-eq-J-symmetry}
    \end{equation}
\end{proposition}
\begin{proof}
    对 \eqref{ch4-eq-Lax-mod-x} 做映射 $ \zeta \to \zeta^{*} $ 有
    \begin{equation}
        J_{x}(\zeta^{*}) = - \rmi \zeta^{*} [\Lambda, J(\zeta^{*})] + Q J(\zeta^{*}),
    \end{equation}
    两边做 Hermitian 转置, 有
    \begin{equation}
        \begin{aligned}
            J_{x}^{\dagger}(\zeta^{*}) = \rmi \zeta [\Lambda, J^{\dagger}(\zeta^{*})] + J^{\dagger}(\zeta^{*}) Q^{\dagger} =  \rmi \zeta [\Lambda, J^{\dagger}(\zeta^{*})] + J^{\dagger}(\zeta^{*}) M_{0} Q M_{0}^{-1} \\
            \implies J_{x}^{\dagger} M_{0} = \rmi \zeta [\Lambda, J^{\dagger}(\zeta^{*}) M_{0}] + J^{\dagger}(\zeta^{*}) M_{0} Q,
        \end{aligned}
    \end{equation}
    与伴随问题 \eqref{ch4-eq-Lax-adjoint} 对比可得 $J^{\dagger}(\zeta^{*}) M_{0} = C J^{-1}(\zeta)$, 其中 $C$ 为常数矩阵, 由边界条件可得 $C = M_{0}$.

    利用 $J_{-} E = J_{+} E S(\zeta)$ 及 $E(\zeta^{*})^{\dagger} = E^{-1}$ 可得
    \begin{equation}
        S^{\dagger}(\zeta^{*}) = [E^{-1}(\zeta^{*}) J_{+}^{-1}(\zeta^{*}) J_{-}(\zeta^{*}) E(\zeta^{*})]^{\dagger} = E^{-1} M_{0} J_{-}^{-1} J_{+} M_{0}^{-1} E = M_{0} S^{-1}(\zeta) M_{0}^{-1}.
    \end{equation}
\end{proof}
进一步可得
\begin{equation}
    [P^{+}(\zeta^{*})]^{\dagger} M_{0} = M_{0} P^{-}(\zeta),
\end{equation}
\begin{proof}
    \begin{equation}
        \begin{aligned}
            (P^{+}(\zeta^{*}))^{\dagger} M_{0} & = (J_{+}(\zeta^{*}) H_{2} + J_{-}(\zeta^{*}) H_{1})^{\dagger} M_{0} = H_{2} J_{+}^{\dagger}(\zeta^{*}) M_{0} + H_{1} J_{-}^{\dagger}(\zeta^{*}) M_{0} \\
                                               & = H_{2} M_{0} J_{+}^{-1}(\zeta) + H_{1} M_{0} J_{-}^{-1}(\zeta) = M_{0}[H_{2} J_{+}^{-1}(\zeta) + H_{1} J_{-}^{-1}(\zeta)] = M_{0} P^{-}(\zeta).
        \end{aligned}
    \end{equation}
\end{proof}
由 \eqref{ch4-eq-RHP-P1} 及 $J_{\pm}$ 定义易见
\begin{equation}
    \det P^{+} = \det(J_{+,1}, J_{+,2}, J_{-,3}) = \det J_{+} \times \begin{vmatrix}
        1 & 0 & s_{13} \\ 0 & 1 & s_{23} \\ 0 & 0 & s_{33}
    \end{vmatrix} = s_{33}.\label{ch4-eq-det-P-s}
\end{equation}
同理有 $\det(P^{-}) = \hat{s}_{33}$, 其中 $\hat{s}_{33} = (S^{-1})_{33}$.
\subsection{RHP及孤子解}
至此我们得到一个 RHP,
\begin{definition}{相关 RH 问题}
    由 $J_{\pm}, J_{\pm}^{-1}$ 的渐进性, 有
    \begin{equation}
        P^{\pm} \to I, \quad \zeta \to \infty, \quad \zeta \in \mathbb{C}_{\pm}
        \label{ch4-eq-RHP-normalization}
    \end{equation}
    在实轴上, 其满足
    \begin{equation}
        P^{-}(x,t,\zeta) P^{+}(x,t,\zeta) = G(x,t,\zeta), \quad \zeta \in \mathbb{R} \label{ch4-eq-RHP-jump}
    \end{equation}
    其中跳跃矩阵
    \begin{equation}
        G = E(H_{1} + H_{2} S(\zeta))(H_{1} + S^{-1}(\zeta) H_{2}) E^{-1}.
    \end{equation}
\end{definition}
\begin{proof}
    由 $P^{\pm}$ 的定义
    \begin{equation}
        \begin{aligned}
            G & = (H_{2} J_{+}^{-1} + H_{1} J_{-}^{-1})(J_{+}H_{2} + J_{-} H_{1})                                                \\
              & = H_{2}^{2} + H_{2} J_{+}^{-1} J_{+} E S E^{-1} H_{1} + H_{1} J_{-}^{-1} J_{-} E S^{-1} E^{-1} H_{2} + H_{1}^{2} \\
              & = H_{2} + H_{1} E S^{-1} E^{-1} H_{2} + H_{2} E S E^{-1} H_{1} + H_{1}                                           \\
              & = E(H_{1} + H_{2} S)(H_{1} + S^{-1} H_{2}) E^{-1}.
        \end{aligned}
    \end{equation}
\end{proof}
接下来只需寻找满足 \eqref{ch4-eq-RHP-normalization} 及 \eqref{ch4-eq-RHP-jump} 的矩阵函数 $ P^{\pm}(\zeta) $.
\subsubsection{正则情形}
若 $\forall \zeta \in \mathbb{C}_{\pm}, \det(P^{\pm}) \neq 0$, 且在实数轴上无奇点, 则为正则 RH 问题, 则由 Plemelj 公式有
\begin{equation}
    (P^{+})^{-1} = I + \frac{1}{2 \pi \rmi} \int_{\mathbb{R}} \frac{(I - G(\lambda)) (P^{+}(\lambda))^{-1}}{\lambda - \zeta} \rmd \lambda, \quad \zeta \in \mathbb{C}_{+}.
\end{equation}
\subsubsection{非正则情形}
若 $\det(P^{\pm})$ 仅存在简单零点, 不妨设在 $\zeta = \zeta_{k}$ 处, $\det(P^{\pm}(\zeta_{k})) = 0$. 为表示方便, 我们记
\begin{equation}
    \{\zeta_{k} \in \mathbb{C}_{+} : s_{33}(\zeta_{k}) = 0\}, \quad
    \{\bar{\zeta}_{k} \in \mathbb{C}_{-} : \hat{s}_{33}(\bar{\zeta}_{k}) = 0\},\quad k = 1, 2, \ldots, N,
\end{equation}
均为简单零点, 则 $\ker(P^{\pm})$ 有唯一的特征向量 $v_{k}, \bar{v}_{k}$, 且满足
\begin{equation}
    P^{+}(\zeta_{k}) v_{k} = 0, \quad \bar{v}_{k} P^{-}(\bar{\zeta}_{k}) = 0, \quad 1 \leq k \leq N.
    \label{ch4-eq-RHP-kernel}
\end{equation}
接下来考虑消去零点结构, 构造 $\Omega := \{(x,t): \det M(x,t) \neq 0\}$, 其中
\begin{equation}
    M_{jk} = \frac{\bar{v}_{j} v_{k}}{\bar{\zeta}_{j} - \zeta_{k}}, \quad 1 \leq j,k \leq N.
\end{equation}
构造单极点 Gamma 矩阵
\begin{equation}
    \Gamma(\zeta) = I + \sum_{j,k = 1}^{N} \frac{v_{j} (M^{-1})_{jk} \bar{v}_{k}}{\zeta - \bar{\zeta}_{k}}, \quad
    \Gamma^{-1}(\zeta) = I - \sum_{j,k = 1}^{N} \frac{v_{j} (M^{-1})_{jk} \bar{v}_{k}}{\zeta - \zeta_{j}}.
    \label{ch4-eq-Gamma-N}
\end{equation}
由 \eqref{ch4-eq-RHP-kernel} 可得
\begin{equation}
    \Gamma(\zeta_{l}) v_{l} = v_{k} + \sum_{j,k}^{N} v_{j} (M^{-1})_{jk} M_{kl} = 0, \quad \bar{v}_{l} \Gamma^{-1}(\bar{\zeta}_{l}) \sum_{j,k}^{N} M_{lj} (M^{-1})_{jk} \bar{v}_{k} = 0, l \in 1, 2, \ldots, N.
    \label{ch4-eq-Gamma-kernel}
\end{equation}
易见 $\Gamma \Gamma^{-1}$ 可能的极点仅为 $\zeta_{l}, \bar{\zeta}_{l}$, 且均为一阶极点. 由留数条件
\begin{equation}
    \Res_{\zeta = \zeta_{l}}(\Gamma \Gamma^{-1}) = - \Gamma v_{l} \sum_{j,k = 1}^{N} (M^{-1})_{jk} \bar{v}_{k} = 0, \quad \Res_{\zeta = \bar{\zeta}_{l}}(\Gamma \Gamma^{-1}) = \sum_{j,k = 1}^{N} v_{j} (M^{-1})_{jk} \bar{v}_{k} \Gamma^{-1}(\bar{\zeta}_{l}) = 0,
\end{equation}
且
\begin{equation}
    \Gamma \Gamma^{-1} \to I, \quad \zeta \to \infty.
    \label{ch4-eq-Gamma-infinity}
\end{equation}
由 Liouville 定理可得 $\Gamma \Gamma^{-1} \equiv I$, 则 $\forall (x,t) \in \Omega, \det(\Gamma) \neq 0$. 由 \eqref{ch4-eq-Gamma-kernel} 可得: $\det \Gamma$ 在 $\zeta_{l}$ 处有简单零点, 在 $\bar{\zeta}_{l}$ 处有简单极点. 结合无穷远处归一化条件有
\begin{equation}
    \det F(\zeta) = \prod_{l=1}^{N} \frac{\zeta - \zeta_{l}}{\zeta - \bar{\zeta}_{l}}.
\end{equation}
定义
\begin{equation}
    \hat{P}^{+} = P^{+} \Gamma^{-1}, \quad \hat{P}^{-} = \Gamma P^{-}, \quad P^{+} = \hat{P}^{+} \Gamma, \quad P^{-} = \Gamma^{-1} \hat{P}^{-}.
    \label{ch4-eq-RHP-decomposition}
\end{equation}
易见它们在 $\zeta_{l}, \bar{\zeta}_{l}$ 处无奇点, 且满足
\begin{equation}
    \Res_{\zeta = \zeta_{l}}\hat{P}^{+} = -P^{+}(\zeta_{l})v_{l} \sum_{k = 1}^{N}(M^{-1})_{lk}\bar{v}_{k} = 0, \quad
    \Res_{\zeta = \bar{\zeta}_{l}}\hat{P}^{-} = \sum_{j = 1}^{N}v_{j}(M^{-1})_{jl}\bar{v}_{l} P^{-}(\bar{\zeta}_{l}) = 0.
\end{equation}
且
\begin{equation}
    \det\hat{P}^{+} = s_{33}(\zeta) \prod_{k = 1}^{N} \frac{\zeta - \bar{\zeta}_{k}}{\zeta - \zeta_{k}},  \quad
    \det\hat{P}^{-} = \hat{s}_{33}(\zeta) \prod_{k = 1}^{N} \frac{\bar{\zeta} - \zeta_{k}}{\bar{\zeta} - \bar{\zeta}_{k}}.
\end{equation}
在各自零点取极限可得
\begin{equation}
    \det \hat{P}^{+}(\zeta_{l}) = s'(\zeta_{l}) (\zeta_{l} - \bar{\zeta}_{l}) \prod_{k \neq l} \frac{\zeta_{l} - \bar{\zeta}_{k}}{\zeta_{l} - \zeta_{k}} \neq 0, \quad
    \det \hat{P}^{-}(\zeta_{l}) = \hat{s}'(\zeta_{l}) (\zeta_{l} - \bar{\zeta}_{l}) \prod_{k \neq l} \frac{\zeta_{l} - \zeta_{k}}{\zeta_{l} - \bar{\zeta}_{k}} \neq 0.
\end{equation}
其中 $s'(\zeta_{l})$ 表示在 $\zeta = \zeta_{l}$ 处的导数. 故 $\hat{P}^{\pm}$ 在 $\mathbb{C}_{\pm}$ 中解析无奇点, 连续到实轴, 且满足 $\hat{P}^{\pm} \to I$. 由 $P^{-} P^{+} = G$ 可得
\begin{equation}
    \hat{P}^{-} \hat{P}^{+} = \Gamma G \Gamma^{-1} = \hat{G}, \quad \zeta \in \mathbb{R}.
\end{equation}
由 \eqref{ch4-eq-Gamma-infinity} 可得 $I - \hat{G} = \Gamma (I - G) \Gamma^{-1}$, 则 $\hat{G} - I \in L^{2}(\mathbb{R}) \cap L^{\infty}(\mathbb{R})$. 由正则 RH 问题的唯一性可得 $\hat{P}^{\pm}$ 唯一存在, 同样有 Plemelj 公式
\begin{equation}
    (\hat{P}^{+})^{-1}(\zeta) = I + \frac{1}{2 \pi \rmi} \int_{\mathbb{R}} \frac{(I - \hat{G}(\lambda)) (\hat{P}^{+}(\lambda))^{-1}}{\lambda - \zeta} \rmd \lambda, \quad \zeta \in \mathbb{C}_{+}.
\end{equation}


































若 RHP \eqref{ch4-eq-RHP-jump} 满足 \eqref{ch4-eq-det-P-s} 及零边界条件, 则可解出所有零点. 注意到 $ P^{+} $ 为 \eqref{ch4-eq-AKNS-Lax-Pair-x} 的解, 利用 Taylor 展开
\begin{equation}
    P^{+}(x,\zeta) = I + \frac{1}{\zeta} P_{1}^{+}(\zeta) + O(\zeta^{-2}), \quad \zeta \to \infty
\end{equation}
将其带入 \eqref{ch4-eq-AKNS-Lax-Pair-x} 比较 $\zeta^{0}$ 的系数可得
\begin{equation}
    Q = \rmi [\Lambda, P_{1}^{+}] = - \begin{pmatrix}
        0 & 0 & 2 \rmi P_{13} \\ 0 & 0 & 2 \rmi P_{23} \\ -2 \rmi P_{31} & -2 \rmi P_{32} & 0
    \end{pmatrix}, \label{ch4-eq-Q-from-P1}
\end{equation}
其中 $ P_{1}^{+} = \left(P_{ij}\right) $,

\subsection{$N$ 孤子解}

为找到孤子解, 考虑无反射势情形, 即令 $G = I $, 有 $\hat{G} = I$. 此时 $\hat{P}^{+} \hat{P}^{-} = I$, 则 $P^{+} = \Gamma$, $P^{-} = \Gamma^{-1}$, 故
\begin{equation}
    P^{+}(\zeta) = I + \sum_{j,k=1}^{N} \frac{v_{j} \left(M^{-1}\right)_{jk} \bar{v}_{k}}{\zeta - \bar{\zeta}_{k}}, \quad M_{jk} = \frac{\bar{v}_{j} v_{k}}{\bar{\zeta}_{j} - \zeta_{k}}
\end{equation}
为找出 $ v_{k}(x,t, \zeta) $, 将 $ P^{+}v_{k} = 0 $ 带入 \eqref{ch4-eq-AKNS-Lax-Pair-x}-\eqref{ch4-eq-AKNS-Lax-Pair-t} 得
\begin{equation}
    \begin{aligned}
        P^{+}_{x} = - \rmi \zeta [\Lambda, P^{+}] + Q P^{+}, \\
        P^{+}_{t} = 2 \rmi \zeta^{2} [\Lambda, P^{+}] + R P^{+},
    \end{aligned}
    \implies
    \begin{aligned}
        P^{+}(\zeta_{k},x) (\frac{d v_{k}}{dx}+ \rmi \zeta_{k} \Lambda v_{k}) = 0, \\
        P^{+}(\zeta_{k},x) (\frac{d v_{k}}{dt} - 2 \rmi \zeta_{k}^{2} \Lambda v_{k}) = 0,
    \end{aligned}
    \implies
    \begin{aligned}
        v_{kx} = - \rmi \zeta_{k} \Lambda v_{k}, \\
        v_{kt} = 2 \rmi \zeta_{k}^{2} \Lambda v_{k},
    \end{aligned}.
\end{equation}
则 $ v_{k}(x,t, \zeta) = E v_{k0}$, 其中 $v_{k0}$ 为常向量, 同理可得 $ \bar{v}_{k}(x,t, \bar{\zeta}) = v_{k0}^{\dagger} M_{0} E^{*} $, 则
\begin{equation}
    p(x,t) = - 2 \rmi P_{13} = - 2 \rmi \left( \sum_{j,k=1}^{N} v_{j}(M^{-1})_{jk} \bar{v}_{k}\right)_{13} \label{ch4-eq-NLS-solution-q}
\end{equation}
其余 $ q, r_{1}, r_{2} $ 类似. 不失一般性, 可取 $ v_{k0} = (\vec{a}, 1) = (\alpha_{k}, \beta_{k}, 1)^{T} $, 则
\begin{equation}
    p(x,t) = - 2 \rmi \sum_{j,k=1}^{N} \alpha_{k} \rme^{-(\theta_{k}^{*}-\theta_{j})}(M^{-1})_{jk},
\end{equation}
不失一般性, 可取 $ v_{k0} = (\vec{a}, 1) = (\alpha_{k}, \beta_{k}, 1)^{T} $, 则
\begin{equation}
    p(x,t) = - 2 \rmi \sum_{j,k=1}^{N} \alpha_{k} \rme^{-(\theta_{k}^{*}-\theta_{j})}(M^{-1})_{jk},
\end{equation}

\section{逆空间约化情形}
或许读者会问, 相较于上一章节 $2 \times 2$ 在原方程做变换, 对 $3 \times 3$ 情形在原方程和伴随问题下做变换, 并直接取等价是否正确? 我们将以逆空间情形为例, 利用传统 RH 方法进行求解, 并最后证明二者的等价性并回答问题 \ref{ch3-question-RH-and-eigenvalue-problem}. 然后在下节直接给出三个定理及简要证明.

\subsection{对称性}
% 实际上, coupled NLS 和 reverse-space NLS 仅对称性不同, 其余内容完全相同, 故我们仅给出对称性证明, 其余内容可参考上一节.
由于仅有耦合 NLS 和逆空间 NLS 仅有对称性不同, 故我们仅给出对称性证明
为表示方便, 我们记 $\phi_{\pm} = J_{\pm} E$, 则 $\phi_{\pm}$ 为 Lax 对 \eqref{ch4-eq-AKNS-Lax-Pair-x} 的解, 且满足对称关系
\begin{equation}
    \phi_{-} = \phi_{+} S(\zeta) \implies J_{-} E = J_{+} E S(\zeta), \quad \zeta \in \mathbb{R},
\end{equation}
则由于 $\det(J_{\pm}) = 1$, 则 $\det(S(\zeta)) = 1$. 且
\begin{align}
    P_{o}^{+} & = (\phi_{+,1}, \phi_{+,2}, \phi_{-,3}) E^{-1} = (J_{+,1}, J_{+,2}, J_{-,3}) = J_{+} H_{2} + J_{-} H_{1}, \quad \zeta \in \mathbb{C}_{+}, \\
    P_{o}^{-} & = (\phi_{-,1}, \phi_{-,2}, \phi_{+,3}) E^{-1} = (J_{-,1}, J_{-,2}, J_{+,3}) = J_{-} H_{2} + J_{+} H_{1}, \quad \zeta \in \mathbb{C}_{-}
\end{align}
这里不同于耦合 NLS, 我们约定记号 $P_{o}^{\pm}$ 表示为原谱问题的列解组成的矩阵. $P_{a}^{\pm}$ 表示为伴随谱问题的行解组成的矩阵.

接下来我们考虑伴随问题
\begin{equation}
    K_{x} = -\rmi \zeta [\Lambda, K] - K Q.
    \label{ch4-eq-reverse-x-Lax-adjoint}
\end{equation}
由原问题及伴随问题关系, 可得 $J^{-1}$ 亦为伴随问题的解, 若设
\begin{equation}
    \phi_{\pm}^{-1} = E^{-1} J_{\pm}^{-1} = (\hat{\phi}_{\pm,1}, \hat{\phi}_{\pm,2}, \hat{\phi}_{\pm,3})^{T},
\end{equation}
重复上述步骤有
\begin{equation}
    \hat{\phi}_{+,1}, \hat{\phi}_{+2}, \hat{\phi}_{-,3} \in \mathbb{C}_{-}, \quad \hat{\phi}_{-,1}, \hat{\phi}_{-,2}, \hat{\phi}_{+,3} \in \mathbb{C}_{+},
\end{equation}
则
\begin{align}
    P_{a}^{+} & = E (\hat{\phi}_{-,1}, \hat{\phi}_{-,2}, \hat{\phi}_{+,3}) = \left([J_{-}^{-1}]_{1}, [J_{-}^{-1}]_{2}, [J_{+}^{-1}]_{3}\right) = H_{2} J_{-}^{-1} + H_{1} J_{+}^{-1}, \quad \zeta \in \mathbb{C}_{+}, \\
    P_{a}^{-} & = E (\hat{\phi}_{+,1}, \hat{\phi}_{+,2}, \hat{\phi}_{-,3}) = \left([J_{+}^{-1}]_{1}, [J_{+}^{-1}]_{2}, [J_{-}^{-1}]_{3}\right) = H_{2} J_{+}^{-1} + H_{1} J_{-}^{-1}, \quad \zeta \in \mathbb{C}_{-}.
\end{align}
易见如下结论:
\begin{proposition}{对称关系}
    逆空间 NLS 方程满足如下对称关系
    \begin{equation}
        M_{1} J_{\pm}^{-1} = J_{\mp}^{\dagger}(-x,t, -\zeta^{*}) M_{1}, \quad S^{\dagger}(-\zeta^{*}) = M_{1} S M_{1}^{-1}
        \label{ch4-eq-reverse-x-symmetry}
    \end{equation}
\end{proposition}
我们约定符号 $(f(-x,t))_{x} = \frac{\rmd f(-x)}{\rmd x}$, 其中 $f_{x}(-x)$ 仅作为符号, 不为求导运算.
\begin{proof}
    对 \eqref{ch4-eq-Lax-mod-x} 两边做 $ x \to -x, \zeta \to -\zeta^{*}$, 有
    \begin{equation}
        \begin{aligned}
            LHS & = (J_{+}(-x, t, -\zeta^{+}))_{x} = \frac{\partial J(g,t,-\zeta^{*})}{\partial g} \frac{\partial g}{\partial x} = -J_{+,x}(-x, t, -\zeta^{*}), \quad g = -x, \\
            RHS & = -\rmi \zeta^{*} [\Lambda, J_{+}(-x,t,-\zeta^{*})] + Q(-x) J_{+}(-x,t,-\zeta^{*}).
        \end{aligned}
    \end{equation}
    两边同时做 Hermitian 转置,
    \begin{equation}
        J_{+,x}^{\dagger}(-x,t,-\zeta^{*}) = -\rmi \zeta [\Lambda, J_{+}^{\dagger}(-x, t, -\zeta^{*})] + J_{+}^{\dagger}(-x,t,-\zeta^{*}) Q^{\dagger}(-x).
    \end{equation}
    由于 $Q^{\dagger}(-x) = M_{1} Q(x) M_{1}^{-1}$, 则
    \begin{equation}
        \begin{aligned}
            J_{+,x}^{\dagger}(-x,t,-\zeta^{*})               & = -\rmi \zeta [\Lambda, J_{+}^{\dagger}(-x, t, -\zeta^{*})] + J_{+}^{\dagger}(-x,t,-\zeta^{*}) M_{1} Q(x) M_{1}^{-1} \\
            \implies J_{+,x}^{\dagger}(-x,t,-\zeta^{*})M_{1} & = -\rmi \zeta [\Lambda, J_{+}^{\dagger}(-x,t,-\zeta^{*})M_{1}] + J_{+}^{\dagger}(-x,t,-\zeta^{*}) M_{1} Q(x),
        \end{aligned}
    \end{equation}
    与 \eqref{ch4-eq-reverse-x-Lax-adjoint} 作比较, 可得 $J_{+}^{\dagger}(-x,t,-\zeta^{*}) M_{1} = C J_{-}^{-1}(x,t,\zeta)$, 其中 $C$ 为常数矩阵, 不失一般性, 取 $C = M_{1}$, 则
    \begin{equation}
        M_{1} J_{+}^{-1}(x,t,\zeta) = J_{-}^{\dagger}(-x,t,-\zeta^{*}) M_{1}.
    \end{equation}
    由 $J_{-}E = J_{+} E S(\zeta)$, 可得
    \begin{equation}
        S^{\dagger}(-\zeta^{*}) = \left(E^{-1}(-x, t, -\zeta^{*}) J_{+}^{-1}(-x, t, -\zeta^{*}) J_{-}(-x, t, -\zeta^{*}) E(-x, t, -\zeta^{*})\right)^{\dagger}.
    \end{equation}
    由 $E^{\dagger}(-x, t, -\zeta^{*}) = E^{-1}$, 则
    \begin{equation}
        \begin{aligned}
            S^{\dagger}(-\zeta^{*}) & = E^{-1} J_{-}^{\dagger}(-x, t, -\zeta^{*}) (J_{+}^{-1}(-x, t, -\zeta^{*}))^{\dagger} E \\
                                    & = E^{-1} M_{1} J_{+}^{-1}(x,t,\zeta) M_{1}^{-1} M_{1} J_{-}(x,t,\zeta) M_{1}^{-1} E     \\
                                    & = M_{1} E^{-1} J_{+}^{-1} J_{-} E M_{1}^{-1}                                            \\
                                    & = M_{1} S(\zeta) M_{1}^{-1}.
        \end{aligned}
    \end{equation}
\end{proof}
利用上面对称关系, 可得
\begin{proposition}{$P_{o}^{\pm}$ 与 $P_{a}^{\pm}$ 对称关系}\label{ch4-prop:reverse-x-P-symmetry}
    \begin{equation}
        M_{1} P_{a}^{\pm}(x, t, \zeta) = \left( P_{o}^{\pm}(-x, t, -\zeta^{*}) \right)^{\dagger} M_{1}, \quad \zeta \in \mathbb{C}_{\pm}.
    \end{equation}
\end{proposition}
\begin{proof}
    \begin{equation}
        \begin{aligned}
            M_{1} P_{a}^{+} & = M_{1} (H_{2} J_{-}^{-1} + H_{1} J_{+}^{-1}) = H_{2} M_{1} J_{-}^{-1} + H_{1} M_{1} J_{+}^{-1} \\
                            & = H_{2} J_{+}^{\dagger}(-x,t,-\zeta^{*})M_{1} + H_{1} J_{-}^{\dagger}(-x,t,-\zeta^{*})M_{1}     \\
                            & = (H_{2} J_{+}^{\dagger}(-x,t,-\zeta^{*}) + H_{1} J_{-}^{\dagger}(-x,t,-\zeta^{*})) M_{1}       \\
                            & = (P_{o}^{+}(-x,t,-\zeta^{*}))^{\dagger} M_{1}.
        \end{aligned}
    \end{equation}
\end{proof}

标准 RH 方法中的 $P^{\pm}$ 与本问题中的 $P_{o}^{\pm}$ 和 $P_{a}^{\pm}$ 有如下对应关系:
\begin{equation}
    P^{+} = P_{o}^{+}, \quad P^{-} = P_{a}^{-},
\end{equation}
为表示方便, 我们接下来仍以 $P^{\pm}$ 表示.

\begin{equation}
    \begin{aligned}
        P_{o}^{+}(\zeta_{k}) v_{k} = 0, \quad \hat{v}_{k} P_{a}^{+}(\hat{\zeta}_{k})  = 0,          \\
        P_{o}^{-}(\hb{\zeta}_{k}) \hb{v}_{k} = 0, \quad \bar{v}_{j} P_{a}^{-}(\bar{\zeta}_{j}) = 0. \\
    \end{aligned}
\end{equation}

\begin{remark}
    注意这里由于 $\hat{\zeta}_{k} = -\zeta_{k}^{*}$, 则我们无法在利用原问题处理 $\mathbb{C}_{+}$ 的同时得到 $\zeta \in \mathbb{C}_{-}$ 的结果, 故只能使用这样的方式求解, 下面得到的特征向量也与标准 RH 问题中不同. 这是其与耦合 NLS 最大的不同之处.
\end{remark}

利用对称关系 \eqref{ch4-eq-reverse-x-symmetry} 可得
\begin{equation}
    (P_{o}^{+}(-x,t,-\zeta^{*}))^{\dagger} = M_{1} P_{a}^{+}(x,t,\zeta) M_{1}^{-1}
\end{equation}
则两边同左乘 $v_{k}^{\dagger}(-x,t,-\zeta^{*})$ 有
\begin{equation}
    v_{k}(-x,t,-\zeta^{*})^{\dagger} (P_{o}^{+}(-x,t,-\zeta^{*}))^{\dagger} = v_{k}^{\dagger}(-x,t,-\zeta^{*}) M_{1} P_{a}^{+}(x,t,\zeta) = 0
\end{equation}
与 $\hat{v}_{k} P_{a}^{+} = 0$ 比较有 $\hat{v}_{k} = v_{k}^{\dagger}(-x,t,-\zeta^{*})M_{1}$. 同理有 $\hb{v}_{k} = M_{1}^{-1} \bar{v}_{k}^{\dagger}(-x,t,-\zeta^{*})$.
\subsection{RH 问题及 N 孤子解}
将 $P^{+}$ 在 $\zeta \to \infty$ 处 Taylor 展开有

\begin{equation}
    P^{+} = I + \frac{1}{\zeta} P_{1}^{+}(\zeta) + \mathcal{O}(\zeta^{-2}), \quad \zeta \to \infty
\end{equation}
将上式与 Lax 对 \eqref{ch4-eq-AKNS-Lax-Pair-x} 比较有
\begin{equation}
    Q = \rmi [\Lambda, P_{1}^{+}] = - \begin{pmatrix}
        0 & 0 & 2 \rmi P_{13} \\ 0 & 0 & 2 \rmi P_{23} \\ -2 \rmi P_{31} & -2 \rmi P_{32} & 0
    \end{pmatrix},
\end{equation}
其中 $P_{1}^{+} = (P_{ij})_{3 \times 3}$, 接下来考虑无反射势情形, 即令 $G = I$, 得到孤子解情形
\begin{equation}
    P^{+}(\zeta) = I + \sum_{j,k = 1}^{N} \frac{v_{j} (M^{-1})_{jk} \bar{v}_{k}}{\zeta - \bar{\zeta}_{k}}
    \label{ch4-corrected-eq-reflectionless-P-plus}
\end{equation}
其中 $M$ 为 $N \times N$ 矩阵, 且
\begin{equation}
    M_{jk} = \frac{\bar{v}_{j} v_{k}}{\bar{\zeta}_{j} - \zeta_{k}}.
\end{equation}
为找到 $v_{k}(x,t)$ 具体形式, 将 $P^{+} v_{k} = 0$ 带入 Lax 对 \eqref{ch4-eq-AKNS-Lax-Pair-x}-\eqref{ch4-eq-AKNS-Lax-Pair-t}, 可得
\begin{equation}
    \begin{aligned}
        (P^{+}v_{k})_{x} = P_{x}^{+} v_{k} + P^{+} v_{kx} = \rmi \zeta_{k} P^{+} \Lambda v_{k} + P^{+} v_{kx} = 0  \\
        \implies P^{+}(\rmi \zeta_{k} \Lambda v_{k} + v_{kx}) = 0 \implies v_{kx} = -\rmi \zeta_{k} \Lambda v_{k}. \\
    \end{aligned}
\end{equation}
同理有 $v_{kt} = 2 \rmi \zeta_{k}^{2} \Lambda v_{k}$, 则 $v_{k} = E v_{k0}$, 其中 $v_{k0} = (\alpha_{k}, \beta_{k}, \gamma_{k})^{T}$ 为常向量.

同理 $\bar{v}_{k} = \bar{v}_{k0} E^{-1}$, 其中 $\bar{v}_{k0} = (\bar{\alpha}_{k}, \bar{\beta}_{k}, \bar{\gamma}_{k})$ 为常向量.

利用命题 \ref{ch4-prop:reverse-x-P-symmetry} 的对称关系有 $\hat{v}_{k} = (v_{k}(-x,t,-\zeta^{*}))^{\dagger} M_{1}$, 则
\begin{equation}
    \hat{v}_{k} = v_{k0}^{\dagger} E^{-1} M_{1} \quad \hb{v}_{k} = M_{1}^{-1} E \bar{v}_{k0}^{\dagger}.
\end{equation}

\begin{remark}
    这里是因为 $[\theta(-x,t,-\zeta^{*})]^{*} = -\theta_{k}$, 和 $M_{1}^\dagger = M_{1}$.
\end{remark}
至此, 我们已使用标准 RH 方法求解逆空间 NLS 方程形式上的 $N$ 孤子解, 但是缺少 $v_{k}, \bar{v}_{k}$ 的具体形式和约束条件, 下面我们考虑这一点.

\subsubsection{特征向量约束关系}\label{ch4-sec-reverse-x-discrete-spectrum-relation}
现在我们考虑同半平面关系标准无反射公式 \eqref{ch4-corrected-eq-reflectionless-P-plus} 中 $P_{\mathrm o}^{+}$ 的列核 $v_{k}$ 与 $P_{\mathrm a}^{-}$ 的行核 $\bar v_{j}$ 的关系.

由于逆空间约化不联系 $\mathbb{C}_{\pm}$, 一般不能将 $\bar v_{j}$ 逐项写成 $v_{j}$ 的某种形式. 离散向量的正确限制应从重构势满足非局域约化这一条件得到, 考虑 \eqref{ch4-corrected-eq-reflectionless-P-plus} 的非对角元
\begin{equation}
    P_{1,off}^{+} = \begin{pmatrix}
        0 & C_{12} \\ C_{21} & 0
    \end{pmatrix}
\end{equation}
则
\begin{equation}
    \begin{aligned}
        C_{12}(x,t) & = \sum_{j,k=1}^{N} (\alpha_{j}, \beta_{j})^{T} (M^{-1})_{jk} \bar{\gamma}_{k}   \\
        C_{21}(x,t) & = \sum_{j,k=1}^{N} \gamma_{j} (M^{-1})_{jk} (\bar{\alpha}_{k}, \bar{\beta}_{k})
    \end{aligned}
\end{equation}
由 $Q = \rmi [\Lambda, P_{1}^{+}]$ 可得势矩阵 $Q$ 的非对角块为
\begin{equation}
    Q_{off} = \begin{pmatrix}
        0 & -2 \rmi C_{12} \\ 2 \rmi C_{21} & 0,
    \end{pmatrix}
\end{equation}
故
\begin{equation}
    (p, q)^{T}= -2 \rmi C_{12}, \quad (r_{1}, r_{2}) = 2 \rmi C_{21}.
\end{equation}
由 $Q^{\dagger}(-x) = M_{1} Q(x) M_{1}^{-1}$ 可得
\begin{equation}
    C_{21}(x,t) = - C_{12}^{\dagger}(-x,t) A.
    \label{ch4-eq-reverse-x-discrete-spectrum-relation}
\end{equation}
将上式联立可得如下关系
\begin{equation}
    \sum_{j,k=1}^{N} \gamma_{j} (M^{-1})_{jk}(x,t) (\bar{\alpha}_{k}, \bar{\beta}_{k}) = - [\sum_{j,k=1}^{N} (\alpha_{j}(-x), \beta_{j}(-x))^{T} (M^{-1})_{jk}^{\dagger}(-x,t) \bar{\gamma}_{k}]^{\dagger} A.
    \label{ch4-eq-reverse-x-discrete-spectrum-relation-1}
\end{equation}
至此我们终于得到同时涉及 $\mathbb{C}_{\pm}$ 的离散谱向量关系.
\subsubsection{单孤子解}
现在我们考虑 $N = 1$ 情形, 为表示方便我们做如下约定
\begin{equation}
    \begin{aligned}
        \zeta_{1} = \rmi \eta_{1} \in \rmi \mathbb{R}_{+}, \quad v_{1} = (
        \vec{a}_{1} e^{-\theta_{1}}, c_{1} e^{\theta_{1}})^{T}                                                                                               \\
        \bar{\zeta}_{1} = \rmi \bar{\eta}_{1} \in \rmi \mathbb{R}_{-}, \quad \bar{v}_{1} = (\vec{b}_{1} e^{\bar{\theta}_{1}} ,  d_{1} e^{-\bar{\theta}_{1}}) \\
        \theta_{1} = \eta_{1} x - 2 \rmi \eta_{1}^{2} t, \quad \bar{\theta}_{1} = \bar{\eta}_{1} x - 2 \rmi \bar{\eta}_{1}^{2} t
    \end{aligned}
\end{equation}
则易见 $\theta(-x)^{*} = -\theta(x), \bar{\theta}(-x)^{*} = -\bar{\theta}(x)$, 其中 $\vec{a}_{1}$ 为 $1\times 2$ 列向量, $\vec{b}_{1}$ 为 $2\times 1$ 行向量.

不失一般性可取 $c_{1} = d_{1} = 1$, 则有
\begin{equation}
    M_{11} = \frac{\bar{v}_{1} v_{1}}{\bar{\zeta}_{1} - \zeta_{1}} = \frac{\vec{b}_{1} \vec{a}_{1} e^{-(\theta_{1} - \bar{\theta}_{1})} + e^{\theta_{1} - \bar{\theta}_{1}}}{\rmi (\bar{\eta}_{1} - \eta_{1})} = \frac{D_{11}}{\rmi (\bar{\eta}_{1} - \eta_{1})}
\end{equation}
其中 $D_{jk}(x) = \rho_{jk} e^{-(\theta_{k} - \bar{\theta}_{j})} +  e^{\theta_{k} - \bar{\theta}_{j}}, \rho_{jk} = \vec{b}_{j} \vec{a}_{k}$, 则可得
\begin{equation}
    \begin{aligned}
        C_{12}(x,t) & = \vec{a}_{1} M_{11}^{-1} = \frac{\rmi (\bar{\eta}_{1} - \eta_{1})}{D_{11}} \vec{a}_{1} e^{-(\theta_{1}+\bar{\theta}_{1})}, \\
        C_{21}(x,t) & = \vec{b}_{1} M_{11}^{-1} = \frac{\rmi (\bar{\eta}_{1} - \eta_{1})}{D_{11}} \vec{b}_{1} e^{\theta_{1}+\bar{\theta}_{1}}.
    \end{aligned}
\end{equation}
将上式带入 \eqref{ch4-eq-reverse-x-discrete-spectrum-relation} 可得
\begin{equation}
    \begin{aligned}
        -C_{12}^{\dagger}(-x)A & = \frac{\rmi(\bar{\eta}_{1} - \eta_{1})}{D_{11}^{*}(-x)}\vec{a}_{1}^{\dagger}e^{\theta_{1}+\bar{\theta}_{1}} A \\
        C_{21}(x)              & = \frac{\rmi (\bar{\eta}_{1} - \eta_{1})}{D_{11}(x)} \vec{b}_{1} e^{\theta_{1}+\bar{\theta}_{1}}
    \end{aligned}
    \implies \frac{\vec{b}_{1}}{D_{11}(x)} = \frac{\vec{a}_{1}^{\dagger}A}{D_{11}^{*}(-x)}.
\end{equation}
记 $z = e^{\theta_{1} - \bar{\theta}_{1}}$, 则 $D_{11}(x) = \rho_{11} z^{-1} + z, D_{11}^{*}(-x) = \rho_{11}^{*} z + z^{-1}$, 则有
\begin{equation}
    \vec{b}_{1} (\rho_{11}^{*} z + z^{-1}) =\vec{a}_{1}^{\dagger}A (\rho_{11} z^{-1} + z)
\end{equation}
比较 $z$ 的次数有
\begin{equation}
    \begin{aligned}
        \vec{b}_{1} \rho_{11}^{*} & =\vec{a}_{1}^{\dagger}A,          \\
        \vec{b}_{1}               & =\vec{a}_{1}^{\dagger}A \rho_{11}
    \end{aligned} \implies |\rho_{11}|^{2} = 1 \implies \rho_{11} = e^{\rmi \omega}, \omega \in \mathbb{R}.
\end{equation}
又由于 $\rho_{11} = \vec{b}_{1}\vec{a}_{1}$, 则有
\begin{equation}
    \rho_{11} = \vec{b}_{1} \vec{a}_{1} =
    \frac{1}{\rho_{11}^{*}} \vec{a}_{1}^{\dagger} A \vec{a}_{1} \implies\vec{a}_{1}^{\dagger}A \vec{a}_{1} = 1.
\end{equation}
则有
\begin{equation}
    \vec{a}_{1}^{\dagger} A \vec{a}_{1} = 1, \quad \vec{b}_{1} = \rho_{11} \vec{a}_{1}^{\dagger} A, \quad |\rho_{11}| = 1
\end{equation}
对于单孤子情形, 只需满足 $\eta = -\bar{\eta}$, 则为有界到边界(无增减性)的孤子.
\subsubsection{多孤子解}
接下来我们考虑平衡谱意义下的多孤子解,
\begin{equation}
    \begin{aligned}
        M_{jk}                        & = \frac{(\vec{b}_{j} \bar{E}(\bar{\zeta}_{j})) E(\zeta_{k}) (\vec{a}_{k}, 1)^{T}}{\bar{\zeta}_{j} - \zeta_{k}} = \frac{\rho_{jk} e^{\theta(\bar{\zeta}_{j}) + \theta(\zeta_{k})} + e^{-( \theta(\bar{\zeta}_{j}) + \theta(\zeta_{k})}}{\bar{\zeta}_{j} - \zeta_{k}}, \\
        M_{jk}^{*}(-x, t, -\zeta^{*}) & = \frac{(\vec{b}_{j}^{*} E^{*}(-x, t, -\zeta_{j}^{*})) (E^{-1}(-x, t, -\zeta_{k}^{*}))^{*} (\vec{a}_{k}^{*}, 1)^{T}}{(-\zeta_{j}^{*} + \zeta_{k}^{*})^{*}}                                                                                                           \\
                                      & = \frac{\rho_{jk}^{*} e^{-(\theta(\bar{\zeta}_{j}) + \theta(\zeta_{k}))} + e^{\theta(\bar{\zeta}_{j}) + \theta(\zeta_{k})}}{-(\bar{\zeta}_{j} - \zeta_{k})}.
    \end{aligned}
\end{equation}
易见若设 $|\rho_{jk}| = 1$, 则有
\begin{equation}
    M_{jk}(-x,t, -\zeta^{*}) = = \frac{e^{\theta(\bar{\zeta}_{j}) + \theta(\zeta_{k})} + e^{-( \theta(\bar{\zeta}_{j}) + \theta(\zeta_{k})}}{-(\bar{\zeta}_{j} - \zeta_{k})} \implies M_{jk} (-x, t, -\zeta^{*}) = - M_{jk}^{*},
\end{equation}
则有
\begin{equation}
    \begin{aligned}
        \begin{pmatrix} p^{*}(-x,t) \\ q^{*}(-x,t)) \end{pmatrix} & = 2\rmi C_{12} = 2 \rmi \sum_{j,k=1}^{N} \left(\vec{a}_{k}(-x,t,-\zeta_{k}^{*}) (M^{-1}(-x,t,-\zeta_{j}, \bar{\zeta}_{k}^{*}))_{jk} e^{\bar{\theta}(x, t, -\bar{zeta}^{*})} \right)^{*} \\
                                                                  & = - 2 \rmi \sum_{j,k=1}^{N} (M^{-1})_{jk} e^{\theta - \bar{\theta}} \begin{pmatrix}
                                                                                                                                            \alpha_{k}^{*} \\ \beta_{k}^{*}
                                                                                                                                        \end{pmatrix}
    \end{aligned}
\end{equation}
再利用 $Q$ 的约化有
\begin{equation}
    (r_{1}, r_{2}) = -(p^{*}(-x,t), q^{*}(-x,t)) A = 2\rmi \sum_{j,k=1}^{N} (M^{-1})_{jk} e^{\theta - \bar{\theta}} (\alpha_{k}^{*}, \beta_{k}^{*}) A
\end{equation}
若设
\begin{equation}
    v_{k0} = v_{0} = (\alpha, \beta, 1)^{T}, \quad \bar{v}_{j0} = \bar{v}_{0} = (\bar{\alpha}, \bar{\beta}, 1), \quad \forall j, k \in 1, 2, \dots N
\end{equation}
则有
\begin{equation}
    \vec{b}_{j0} = \vec{a}_{j0}^{\dagger} A \implies \bar{v}_{k0} = v_{k0}^{\dagger} M_{0}.
\end{equation}
且满足 $\bar{v}_{0} v_{0} = 2$.
\subsection{特征函数对称性与 RH 方法等价性}
接下来我们证明利用特征函数对称性得到的特征值和特征向量与 RH 方法得到的结果等价, 也即证明下一节中定理的正确性.

考虑特征值和伴随特征值问题
\begin{equation}
    \begin{aligned}
        Y_{x} & = -\rmi \zeta \Lambda Y + Q Y, \\
        K_{x} & = \rmi \zeta K \Lambda - K Q.
    \end{aligned}
\end{equation}
的特征向量分别为 $v_{k}'$ 和 $\hat{v}_{k}'$. 不失一般性, 我们仅考虑驻定情形.
\begin{equation}
    \begin{aligned}
        P^{+}(\zeta_{k}) v_{k} & = (\phi_{+,1}, \phi_{+,2}, \phi_{-,3}) E^{-1} v_{k} = (\phi_{+,1}, \phi_{+,2}, \phi_{-,3}) E^{-1} E v_{k0} = 0, \\
                               & = \alpha_{k} \phi_{+,1} + \beta_{k} \phi_{+,2} + \gamma_{k} \phi_{-,3} = 0.
    \end{aligned}
    \label{ch4-eq-reverse-x-phi-linear-combination}
\end{equation}
其中
\begin{equation}
    \begin{aligned}
        \phi_{+,1} \sim (e^{\rmi \zeta_{k} x}, 0,0)^{T}, \quad \phi_{+,2} \sim (0, e^{\rmi \zeta_{k} x}, 0)^{T}, \quad x \to +\infty \\
        \phi_{-,3} \sim (0,0, e^{- \rmi \zeta_{k} x})^{T}, \quad x \to -\infty
    \end{aligned}
\end{equation}
又由于
\begin{equation}
    Y_{kx} \sim -\rmi \zeta_{k} \Lambda Y_{k} \implies Y_{k} = e^{-\rmi \zeta_{k} x \Lambda} c_{\pm} \quad x \to \pm \infty
\end{equation}
其中 $c_{\pm}$ 为常向量, 利用离散谱条件有
\begin{equation}
    Y_{k} \to (\alpha_{k}' e^{\rmi \zeta_{k} x}, \beta_{k}' e^{\rmi \zeta_{k} x},0)^{T}, \quad x \to +\infty \\
    Y_{k} \to (0,0, -\gamma_{k}' e^{- \rmi \zeta_{k} x})^{T}, \quad x \to -\infty.
\end{equation}
可设 $c_{+} = (\alpha_{k}', \beta_{k}', 0)^{T}, c_{-} = (0,0, - \gamma_{k}')^{T}$, 则在 $x \in \mathbb{R}$ 上有
\begin{equation}
    Y_{k} = \alpha_{k}' \phi_{+,1} + \beta_{k}' \phi_{+,2} = -\gamma_{k}' \phi_{-,3}
\end{equation}
则有
\begin{equation}
    \alpha_{k}' \phi_{+,1} + \beta_{k}' \phi_{+,2} + \gamma_{k}' \phi_{-,3} = 0
    \label{ch4-eq-reverse-x-phi-linear-combination-2}
\end{equation}
故 \eqref{ch4-eq-reverse-x-phi-linear-combination} 与 \eqref{ch4-eq-reverse-x-phi-linear-combination-2} 等价, 故 $ v_{k0} = c v_{k0}'$, 不失一般性, 可令 $c=1$, 则 $\alpha_{k} = \alpha_{k}', \beta_{k} = \beta_{k}', \gamma_{k} = \gamma_{k}'$, 也即 RH 方法得到的孤子解与特征函数对称性得到的孤子解等价.

现在我们可以对问题 \ref{ch3-question-RH-and-eigenvalue-problem} 做肯定的回答. 且不但可以对原问题求对称性, 也可联立原问题和伴随问题求解对称性.

\section{三种约化下的散射数据对称性}
本节我们仿照上一章节, 先考虑 \eqref{ch4-eq-reduction-reverse-x}-\eqref{ch4-eq-reduction-reverse-x-t} 三种约化下的散射数据对称性. 下一节考虑增加二重约化 \eqref{ch4-eq-reduction-q1}-\eqref{ch4-eq-reduction-q3} 时的散射数据对称性.

\subsection{三种约化对称性}
为表示方便, 我们约定符号 $ \zeta_{k} \in \mathbb{C}_{+}, \bar{\zeta}_{k} \in \mathbb{C}_{-} $ 为离散特征值, $ v_{k0}, \bar{v}_{k0} $ 为对应的特征向量. $\hat{\zeta}, \hat{\bar{\zeta}}$ 为我们通过对称关系找到的特征值
\iffalse
    \begin{theorem}[逆空间] %\label{ch4-thm:inverse-space}
        对于逆空间 NLS 方程 \eqref{ch4-eq-reduction-inverse-x-p}-\eqref{ch4-eq-reduction-inverse-x-q}, 如果 $ \zeta $ 是一个特征值, 则 $ -\zeta^{*} $ 也是一个特征值. 因此，非纯虚特征值总是以 $ (\zeta, -\zeta^{*}) $ 的形式成对出现，它们位于复平面的同一半平面。对称关系如下:
        \begin{enumerate}
            \item 若 $ (\zeta_{k}, \hat{\zeta_{k}}) \in \mathbb{C_{+}} $,对应的特征向量为 $ v_{k0}, \hat{v}_{k0}$, 则满足 $\hat{v}_{k0} = {v}_{k0}^{\dagger}M_{1}$.
            \item
                  若 $ (\bar{\zeta_{k}}, \hb{\zeta_{k}}) \in \mathbb{C_{-}} $, 对应的特征向量为 $ \bar{v}_{k0}, \hb{v}_{k0}$, 则满足满足 $\hb{v}_{k0} = \bar{v}_{k0}^{*}M_{1}^{T}$.
        \end{enumerate}
        更进一步, 若 $\zeta_{k} \in i\mathbb{R}_{\pm}$, 则 $\hat{\zeta}_{k} = -\zeta_{k}^{*} = \zeta_{k}$, 我们设 $ a = d $, 且 $ b \in \mathbb{R} $ 为实数, 则
        \begin{enumerate}
            \item 若 $\zeta_{k} \in \mathbb{R}_{+}$ 为纯虚数, 则对应的特征向量为 $v_{k0} = (\alpha_{k}, \alpha_{k}^{*}, 1)^{T}$.
            \item 若 $\zeta_{k} \in \mathbb{R}_{-}$ 为纯虚数, 则对应的特征向量为 $\bar{v}_{k0} = (\bar{\alpha}_{k}, \bar{\alpha}_{k}^{*}, 1)^{T}$.
        \end{enumerate}
    \end{theorem}
\fi
\begin{theorem}[逆空间] \label{ch4-thm:inverse-space}
    对于逆空间 NLS 方程\eqref{ch4-eq-reduction-inverse-x-p}-\eqref{ch4-eq-reduction-inverse-x-q}, 如果 $ \zeta $ 是一个特征值, 则 $ -\zeta^{*} $ 也是一个特征值. 因此，非纯虚特征值总是以 $ (\zeta, -\zeta^{*}) $ 的形式成对出现，它们位于复平面的同一半平面。对称关系如下:
    \begin{equation}
        C_{21}(x,t) = - C_{12}^{\dagger}(-x,t) A.
    \end{equation}
    特别的, 若设
    \begin{equation}
        v_{j0} = v_{0} = (\alpha, \beta, 1)^{T}, \quad \bar{v}_{j0} = \bar{v}_{0} = (\bar{\alpha}, \bar{\beta}, 1), \quad \forall j, k \in 1, 2, \dots N
    \end{equation}
    则满足 $\bar{v}_{j0} = v_{k0}^{\dagger} M_{0}$, 且 $\bar{v}_{j0} v_{j0} = 2$.

    进一步当上下半平面各包含一个纯虚特征值时. 必有
    \begin{equation}
        \zeta_{k} = \rmi \eta_{k} \in \rmi \mathbb{R}_{+}, \quad \bar{\zeta}_{k} = -\rmi \bar{\eta}_{k} \in \rmi \mathbb{R}_{-}, \quad v_{k0} = (\vec{a}_{k}, 1)^{T}, \quad \bar{v}_{k0} = (\vec{b}_{k}, 1),
    \end{equation}
    满足
    \begin{equation}
        \vec{b}_{k} = \rho_{11} \vec{a}_{k}^{\dagger} A_{1}, \quad \vec{a}_{k}^{\dagger} A_{1} \vec{a}_{k} = 1, \quad 1 \leq i, j \leq N
    \end{equation}
\end{theorem}
在证明开始之前, 我们仿照上一章节, 先建立离散数据之间的关系, 考虑离散数据 $ \{\zeta_{k}, \bar{\zeta}_{k}, \alpha_{k}, \beta_{k}, \bar{\alpha}_{k}, \bar{\beta}_{k}, 1 \leq k \leq N\} $ 及原谱问题 $Y_{x} = MY$ 和伴随谱问题 $K_{x} = -KM $
\begin{align}
    Y_{x} & = - \rmi \zeta \Lambda Y + Q Y \label{eq:3-NLS-eigenvalue-problem},       \\
    K_{x} & = \rmi \zeta K \Lambda - K Q \label{eq:3-NLS-adjoint-eigenvalue-problem},
\end{align}
考虑驻定情形, 其中
\begin{equation}
    Q(x)=\begin{pmatrix}
        0           & 0           & p(x,0) \\
        0           & 0           & q(x,0) \\
        -r_{1}(x,0) & -r_{2}(x,0) & 0
    \end{pmatrix}
\end{equation}
这里 $ p(x,0), q(x,0), r_{1}(x,0), r_{2}(x,0) $ 分别为 $ p(x,t), q(x,t), r_{1}(x,t), r_{2}(x,t) $ 在 $ t=0 $ 时的初值. 离散散射数据与离散特征值及特征向量的关系如下: 若 $ \zeta_{k} \in \mathbb{C}_{+} $ 为原谱问题 \eqref{eq:3-NLS-eigenvalue-problem} 的离散特征值, 则存在非零向量 $ v_{k0} = (\alpha_{k}, \beta_{k}, \gamma_{k})^{T} $, 使得对应的离散特征函数 $ Y_{k}(x) $ 满足
\begin{align}
    Y_{k}(x) & \to \left(\alpha_{k} \rme^{\rmi \zeta_{k} x}, \beta_{k} \rme^{\rmi \zeta_{k} x}, 0\right)^T, \quad x \to + \infty,                                  \\
    \quad
    Y_{k}(x) & \to \left(0, 0, -\gamma_{k} \rme^{-\rmi \zeta_{k} x}\right)^T, \quad x \to -\infty \label{eq:large-x-asymptotics-of-reverse-x-3d-NLS-eigenfunction}
\end{align}
同理, 若 $ \bar{\zeta}_{k} \in \mathbb{C}_{-} $ 为伴随谱问题 \eqref{eq:3-NLS-adjoint-eigenvalue-problem} 的离散特征值, 则存在非零向量 $ \bar{v}_{k0} = (\bar{\alpha}_{k}, \bar{\beta}_{k}, \bar{\gamma}_{k}) $, 使得对应的离散特征函数 $ K_{k}(x) $ 满足
\begin{align}
    K_{k}(x) & \to \left(\bar{\alpha}_{k} \rme^{-\rmi \bar{\zeta}_{k} x},\bar{\beta}_{k} \rme^{-\rmi \bar{\zeta}_{k} x}, 0\right)M_{0}, \quad x \to +\infty,                           \\
    K_{k}(x) & \to \left(0, 0, -\bar{\gamma}_{k}\rme^{\rmi \bar{\zeta}_{k}x}\right)M_{0}, \quad x \to -\infty \label{eq:large-x-asymptotics-of-reverse-x-3d-NLS-adjoint-eigenfunction}
\end{align}
鉴于此联系, 为了推导(离散)散射数据的对称关系, 我们将利用特征值问题中离散特征模的对称关系 \eqref{eq:3-NLS-eigenvalue-problem}-\eqref{eq:3-NLS-adjoint-eigenvalue-problem}.
\begin{proof}
    此时 $ r_{1}(x,t) = a p^{*}(-x,t) + b^{*} q^{*}(-x,t), r_{2}(x,t) = b p^{*}(-x,t) + d q^{*}(-x,t) $, 容易发现 $ Q^{\dagger}(-x) = M_{1} Q(x) M_{1}^{-1} $, 所以我们对 \eqref{eq:3-NLS-eigenvalue-problem} 式先共轭转置, 再进行变量替换 $ x \to -x $, 则有
    \begin{equation}
        \begin{aligned}
            Y_{x} = - \rmi \zeta \Lambda Y + Q Y \implies -Y_{x}(-x) & = - \rmi \zeta \Lambda Y(-x) + Q(-x) Y(-x)                                      \\
            \implies -Y_{x}^{\dagger}(-x)                            & = \rmi \zeta^{*} Y^{\dagger}(-x) \Lambda + Y^{\dagger}(-x) Q^{\dagger}(-x)      \\
                                                                     & =\rmi \zeta^{*} Y^{\dagger}(-x)\Lambda + Y^{\dagger}(-x) M_{1} Q(x) M_{1}^{-1},
        \end{aligned}
    \end{equation}
    故
    \begin{equation}
        \hat{Y}_{x}(x) = \hat{Y}(x) [\rmi \hat{\zeta} \Lambda - Q(x)] \label{ch4-eq-reverse-x-symmetry-relation}
    \end{equation}
    其中
    \begin{equation}
        \hat{Y}(x) = Y(-x)^{\dagger}M, \quad \hat{\zeta} = -\zeta^{*}, \label{ch4-eq-reverse-x-symmetry-parameters}
    \end{equation}
    故 \eqref{ch4-eq-reverse-x-symmetry-relation} 显示若 $ \zeta_{k} \in \mathbb{C}_{+} $ 为散射问题 \eqref{eq:3-NLS-eigenvalue-problem} 的特征值, 则 $ \hat{\zeta}_{k} = -\zeta_{k}^{*} \in \mathbb{C}_{+} $ 亦为特征值. 此外, $\zeta_{k}$ 的特征函数 $ Y_{k}(x) $ 与 $\zeta_{k}$ 的特征函数 $ Y_{k}(x) $ 之间满足关系式 \eqref{ch4-eq-reverse-x-symmetry-parameters}.
    但是注意到 $(\zeta, -\zeta^{*})$ 属于同半平面, 故未能联系起 $P^{\pm}$及对应特征向量. 我们需要通过约化关系及 $P^{\pm}$ 定义求解, 具体证明过程见 \ref{ch4-sec-reverse-x-discrete-spectrum-relation} 小节.
\end{proof}
\begin{theorem}[逆时间] \label{ch4-thm:inverse-time}
    这种情况下 $ b = c $, $ p_{1}=p(x,-t), q_{1} = q(x,-t) $

    对于逆时间 NLS 方程 \eqref{ch4-eq-reduction-inverse-t-p}-\eqref{ch4-eq-reduction-inverse-t-q}. 若 $ \zeta $ 是特征值, 则 $ -\zeta $ 也是特征值. 因此, 离散谱关于原点对称, 特征值总是以 $(\zeta, -\zeta)$ 的形式成对出现 ，位于复平面的相对半平面中.

    若 $ \zeta_{k} \in \mathbb{C}_{+} $ 且 $ \bar{\zeta}_{k} = -\zeta_{k} \in \mathbb{C}_{-} $, 其对应的特征向量 $ v_{k0} $ 和 $ \bar{v}_{k0} $ 满足关系 $ \bar{v}_{k0} = v_{k0}^{T}M_{2} $.
\end{theorem}
\begin{proof}
    显然有 $ Q^{T}(x,-t) = -M_{2} Q(x,-t) M_{2}^{-1} $, 对 \eqref{eq:3-NLS-eigenvalue-problem} 同时转置, 对 Lax 对做替换 $t \to -t$ 得
    \begin{equation}
        \begin{aligned}
            Y_{x}(x,-t)                    & = \rmi \zeta \Lambda Y(x,-t) + Q(x,-t)Y(x,-t)                                                                                     \\
            \implies Y_{x}^{T}(x,-t)       & = \rmi \zeta Y^{T}(x,-t) \Lambda + Y^{T}(x,-t) Q^{T}(x,-t) = \rmi \zeta Y^{T}(x,-t) \Lambda - Y^{T}(x,-t) M_{2} Q(x,t) M_{2}^{-1} \\
            \implies Y_{x}^{T}(x,-t) M_{2} & = \rmi \zeta Y^{T}(x,-t) \Lambda M_{2} - Y^{T}(x,-t) M_{2} Q(x,-t)                                                                \\
            \implies \hat{Y}_{x}(x,t)      & = -\hat{Y}(x,t)\left[\rmi \hat{\zeta}\Lambda + Q(t)\right],
        \end{aligned}
    \end{equation}
    其中
    \begin{equation}
        \hat{Y}(t) = Y(-t)^{T}M_{2}, \quad \hat{\zeta} = -\zeta
    \end{equation}
    这意味着 $ (\hat{\zeta}, \hat{Y}(t)) $ 满足伴随特征值问题. 因此, 若 $ \zeta_{k} \in \mathbb{C_{+}} $ 为原问题的特征值, 则 $ \hat{\zeta}_{k} = -\zeta_{k} \in \mathbb{C}_{-} $ 亦为伴随散射问题的特征值, 故 $ \hat{v}_{k} = v_{k}^{T} M_{2} $.
\end{proof}
\begin{remark}
    因为逆时间得到了上下半平面的对称关系, 所以我们不需要通过全局对称性(约化关系)进行推导, 而逆空间和逆时间仅得到了同半平面的对称关系, 所以需要通过全局对称性得到 $\bar{v}_{k0}$. 因逆时间与上一章节相同, 在此略过具体分析.
\end{remark}

\begin{theorem}[逆时空] \label{ch4-thm:inverse-space-time}
    这种情形下 $ b = c $, $ p_{1} = p(-x,-t), q_{1} = q(-x,-t) $

    对于逆时空方程 \eqref{ch4-eq-reduction-inverse-xt-p}-\eqref{ch4-eq-reduction-inverse-xt-q}, $ \zeta \in \mathbb{C}_{+}$ 可以为任意值, 且 $ \bar{\zeta} \in \mathbb{C}_{-}$ 亦为任意特征值. 因此, 离散谱无对称性限制, 特征值可以位于复平面的任意位置, 然而其对应的特征向量必须满足
    \begin{equation}
        C_{21}(x,t) = C_{12}^{T}(-x, -t) A_{2}
    \end{equation}
    对任意正整数 $N$, 若 $ \zeta_{k} \in \mathbb{C}_{+}, \bar{\zeta}_{k} \in \mathbb{C}_{-} $, 若设
    \begin{equation}
        v_{k0} = v_{0} = (\alpha, \beta, 1)^{T}, \quad \bar{v}_{j0} = \bar{v}_{0} = (\bar{\alpha}, \bar{\beta}, 1), \quad \forall j, k \in 1, 2, \dots N
    \end{equation}
    则其对应的特征向量满足
    \begin{equation}
        \bar{v}_{j0} = \rho_{jk} v_{k0}^{T}M_{2}, \quad \rho v_{j0}^{T} M_{2} v_{k0} = 1+ \rho \label{eq:3-NLS-reverse-x-t-eigenvector-symmetry}
    \end{equation}
    其中 $\rho = \bar{v}_{0} v_{0} = \pm 1$ 为常数.
\end{theorem}
\begin{proof}
    对于逆时空方程情形, 易见 $ Q^{T}(-x) = M_{3} Q(x) M_{3}^{-1} $, 则
    \begin{equation}
        \begin{aligned}
            Y_{x} = - \rmi \zeta \Lambda Y + Q Y & \implies -Y_{x}(-x,-t) = - \rmi \zeta \Lambda Y(x,-t) + Q(-x,-t) Y(-x,-t)                           \\
                                                 & \implies -Y_{x}^{T}(-x,-t) = -\rmi \zeta Y^{T}(-x,-t) \Lambda + Y^{T}(-x,-t) Q^{T}(-x,-t)           \\
                                                 & \implies Y_{x}^{T}(-x,-t) = \rmi \zeta Y^{T}(-x,-t) \Lambda - Y^{T}(-x,-t) M_{3} Q(x,-t) M_{3}^{-1} \\
                                                 & \implies Y_{x}^{T}(-x,-t) M_{3} =  Y^{T}(-x,-t)M_{3} [\rmi \zeta \Lambda - Q(x,-t)]                 \\
                                                 & \implies \hat{Y}_{x}(x,t) = \hat{Y}(x,t) [\rmi \hat{\zeta} \Lambda - Q(x,t)]
        \end{aligned}
    \end{equation}
    其中
    \begin{equation}
        \hat{Y}(x,t) = Y(-x,-t)^{T} M_{3}, \quad \hat{\zeta} = \zeta
    \end{equation}
    这意味着任意 $ \zeta_{k} \in \mathbb{C}_{+} $ 可任意取定, 且考虑特征值有
    \begin{equation}
        \begin{aligned}
            M_{jk}        & = \frac{\rho_{jk} e^{\bar{\theta}_{j} + \theta_{k}} + e^{-(\bar{\theta}_{j} + \theta_{k})}}{\bar{\zeta}_{j} - \zeta_{k}} \\
            M_{jk}(-x,-t) & = \frac{\rho_{jk} e^{-(\bar{\theta}_{j} + \theta_{k})} + e^{\bar{\theta}_{j} + \theta_{k}}}{\bar{\zeta}_{j} - \zeta_{k}}
        \end{aligned}
    \end{equation}
    故若 $\rho_{jk} = \bar{v}_{j0} v_{k0} = \pm 1$, 则有 $M_{jk}(-x,-t) = \rho_{jk} M_{jk}$. 则唯有 $\rho_{jk} = 1$ 时, 才能保证 $M_{jk}(-x,-t) = \rho_{jk} M_{jk}$, $\rho_{jk} = -1$ 时, 则 $M_{jk}(-x,-t) = - M_{jk}$.
    \begin{equation}
        (p(-x,-t), q(-x,-t))^{T} = -2 \rmi C_{12} = -2 \rmi \sum_{j,k=1}^{N} \rho_{jk} e^{\theta_{k} + \bar{\theta}_{j}} (M^{-1})_{jk} (\alpha_{k}, \beta_{k})^{T}
    \end{equation}
    利用 $Q$ 的约化关系有
    \begin{equation}
        (r_{1}, r_{2}) = - (p(-x,-t), q(-x,-t))^{T} A_{2} = 2 \rmi \sum_{j,k=1}^{N} \rho_{jk} e^{\theta_{k} + \bar{\theta}_{j}} (M^{-1})_{jk} (\alpha_{k}, \beta_{k})^{T} A_{2}
    \end{equation}
    则若设定
    \begin{equation}
        v_{k0} = v_{0} = (\alpha, \beta, 1)^{T}, \quad \bar{v}_{j0} = \bar{v}_{0} = (\bar{\alpha}, \bar{\beta}, 1), \quad \forall j, k \in 1, 2, \dots N
    \end{equation}
    此时不妨记 $\rho_{jk} = \rho = \pm 1$, 则 $\bar{v}_{j0} = \rho v_{k0}^{T} M_{2}$.
\end{proof}
这一部分的三个定理显示: 相比于 $ 2 \times 2 $ NLS 方程, $ 3\times 3 $ NLS 没有那么好的性质, 只能使用转置和共轭转置的方式求解, 尤其是逆空间和逆时空情形相比 $ 2 \times 2 $ NLS 方程的性质更差, 这是因为二阶矩阵存在 Pauli 矩阵基底, 而三阶矩阵不存在类似的完备基底.
\subsection{三种约化下解的动力学分析}
逆时空作为限制最严格得到的解, 因此也具有最好的性质, 可移动, 可坍缩, 可以展现出大部分局部方程的性质, 因此本小节我们讨论一下它的性质,
\subsubsection{逆时空}
考虑函数
\begin{equation}
    \begin{aligned}
        p(x,t) & = \frac{2 \rmi \alpha_{1} (\zeta_{1} - \bar{\zeta}_{1})}{e^{-2 \rmi \zeta_{1} x + 4 \rmi \zeta_{1}^{2} t} + 2 \alpha_{1} \bar{\alpha}_{1} e^{-2 \rmi \bar{\zeta}_{1} x + 4 \rmi \bar{\zeta}_{1}^{2} t}}                               \\
               & = 2 \frac{\rmi \alpha_{1} (\zeta_{1} - \bar{\zeta}_{1}) e^{2\rmi \zeta_{1} x - 4 \rmi \bar{\zeta}_{1}^{2}t}}{1+ 2\alpha \bar{\alpha}_{1} e^{2 \rmi (\zeta_{1} - \bar{\zeta}_{1}) x - 4 \rmi (\zeta_{1}^{2} - \bar{\zeta}_{1}^{2}) t}}
    \end{aligned}
\end{equation}
则可得波速 $ V = \frac{\Im(\bar{\zeta}_{1}^{2} - \zeta_{1}^{2})}{\Im(\bar{\zeta}_{1} - \zeta_{1})} $, 则
\begin{equation}
    | p | = 2 |\alpha_{1} | |\zeta_{1} - \bar{\zeta}_{1} | \frac{e^{\omega_{1} t}}{1+2\alpha_{1}\bar{\alpha}_{1}e^{\rmi \omega_{2} t}}
\end{equation}
其中
\begin{align}
    \omega_{1} & = 2V \Im(\zeta_{1}) + 4 \Im(\zeta_{1}^{2}),                                         \\
    \omega_{2} & = -2V \Re(\bar{\zeta}_{1} - \zeta_{1}) + 4 \Re(\bar{\zeta}_{1}^{2} - \zeta_{1}^{2})
\end{align}
则若 $ \Re(\omega_{1}) \neq 0 $, 则解在 $ t \to \infty $ 时爆炸; 若 $ \Re(\omega_{2}) \neq 0 $, 则解会周期性的坍塌.
\subsection{原因}
占位
\section{进一步考虑二重约化情形}
本节考虑 $ p, q $ 的约化关系, 占位,